%!TEX root = ../thesis.tex

\chapter{Background and Literature Context}

\section{Compressible Euler Equations}
% <<SECTION_1_BEGIN>>
The Report 1 validation cases use the compressible Euler equations for an
ideal gas. These equations form a suitable first system because they contain
the wave features that a shock-capturing code must resolve, including shocks,
contact discontinuities, and rarefactions, while retaining well documented
Riemann-problem and two-dimensional benchmark references
\citep{toro2009,liska_wendroff_2003}. Magnetic-field effects are therefore
left as project context until the Euler precision and hardware comparisons
have an auditable baseline.

In two space dimensions the conservative form is
\begin{equation}
  \partial_t U + \partial_x F(U) + \partial_y G(U) = 0,
  \qquad
  U =
  \begin{pmatrix}
    \rho\\ \rho u\\ \rho v\\ E
  \end{pmatrix},
\end{equation}
where \(\rho\) is density, \(\rho u\) and \(\rho v\) are the Cartesian momentum
densities, and \(E\) is the total energy density. The physical fluxes are
\begin{equation}
  F(U)=
  \begin{pmatrix}
    \rho u\\ \rho u^2+p\\ \rho uv\\ u(E+p)
  \end{pmatrix},
  \qquad
  G(U)=
  \begin{pmatrix}
    \rho v\\ \rho uv\\ \rho v^2+p\\ v(E+p)
  \end{pmatrix}.
\end{equation}
The one-dimensional form used for shock-tube tests is obtained by suppressing
the \(y\)-dependence and the transverse flux. The ideal-gas closure is
\begin{equation}
  E = \rho e + \frac{1}{2}\rho\left(u^2+v^2\right),
  \qquad
  p = (\gamma-1)\rho e
    = (\gamma-1)\left[E-\frac{1}{2}\rho\left(u^2+v^2\right)\right],
\end{equation}
where \(e\) is the specific internal energy, \(p\) is pressure, and
\(\gamma\) is the ratio of specific heats. This section fixes only the
governing variables and closure; the numerical flux construction is introduced
later with the finite-volume method.
% <<SECTION_1_END>>

\section{Ideal-MHD Project Context}
% <<SECTION_2_BEGIN>>
The longer project extends the Euler model by coupling the gas to a magnetic
field \(\mathbf{B}\). This adds magnetic pressure, magnetic tension, and
magnetosonic wave families to the gas-dynamic system. In nondimensional units
with magnetic permeability set to one, a compact ideal-MHD conservation form is
\begin{equation}
  \partial_t
  \begin{pmatrix}
    \rho\\ \rho\mathbf{u}\\ E\\ \mathbf{B}
  \end{pmatrix}
  +
  \nabla\cdot
  \begin{pmatrix}
    \rho\mathbf{u}\\
    \rho\mathbf{u}\otimes\mathbf{u}
      + p_t I - \mathbf{B}\otimes\mathbf{B}\\
    (E+p_t)\mathbf{u}-(\mathbf{u}\cdot\mathbf{B})\mathbf{B}\\
    \mathbf{u}\otimes\mathbf{B}-\mathbf{B}\otimes\mathbf{u}
  \end{pmatrix}
  =0,
  \qquad
  \nabla\cdot\mathbf{B}=0,
\end{equation}
where \(\mathbf{u}\) is velocity, \(E\) is the total gas, kinetic, and magnetic
energy density, and \(p_t=p+\lvert\mathbf{B}\rvert^2/2\) is the total pressure.
The final flux row is the conservative induction equation for \(\mathbf{B}\).
The constraint \(\nabla\cdot\mathbf{B}=0\) states that the magnetic field is
solenoidal. If numerical divergence is allowed to grow, the computed magnetic
force and MHD wave propagation can be polluted by non-physical monopole error.
Report 1 does not implement or validate MHD; for Report~2, the candidate
numerical choices fall into two families surveyed by \citet{toth_2000}:
hyperbolic divergence cleaning of Dedner-type \citep{dedner_2002}, which
augments the induction equation with a scalar transport that carries
divergence error away, and constrained transport \citep{evans_hawley_1988},
which preserves a discrete staggered-mesh divergence constraint by
construction. The corresponding MHD Riemann solver candidate is the
multi-state HLLD scheme \citep{miyoshi_kusano_2005}, which retains the fast
and Alfv\'en waves missing from a direct Euler-HLLC reuse. A GPU-accelerated
MUSCL--Hancock ideal-MHD precedent is provided by \citet{bard_dorelli_2014}.
Eight-wave formulations \citep{powell_1999} provide an alternative MHD HRSC framework with built-in cleaning, noted here for the Report~2 MHD extension.
% <<SECTION_2_END>>

\section{HRSC and Benchmark Literature}
% <<SECTION_3_BEGIN>>
High-resolution shock-capturing (HRSC) finite-volume methods are a natural
choice for compressible Euler flows because they evolve cell averages in
conservative form while allowing discontinuities to be represented without
tracking shock or contact positions explicitly. In Godunov-type methods, the
face flux is obtained from a local Riemann problem, so the approximate Riemann
solver is not a secondary implementation detail: it controls the dissipation
and wave information used at shocks and contact discontinuities
\citep{toro2009,harten_lax_vanleer_1983}. The Harten--Lax--van Leer (HLL)
family gives a standard approximate-flux framework, and the HLLC extension of
\citet{toro_spruce_speares_1994} restores a contact wave within that family;
the algebraic form used in this report is deferred to Chapter~3.

Benchmark problems serve a different role from the numerical method itself.
They provide recognisable wave patterns and comparison points against which a
solver can be checked before precision and hardware effects are interpreted.
Sod's shock tube and Toro's one-dimensional Riemann problems are standard
sources for shocks, contacts, and rarefactions in a controlled setting
\citep{sod_1978,toro2009}. Liska and Wendroff's two-dimensional Euler Riemann
configurations extend this role to interacting waves and multidimensional
structure without requiring the background chapter to catalogue individual test
cases \citep{liska_wendroff_2003}.
% <<SECTION_3_END>>

\section{Floating-Point Arithmetic and Reproducibility}
% <<SECTION_4_BEGIN>>
IEEE floating-point arithmetic uses finite significand and exponent fields, so
the update is not evaluated over the real numbers. In round-to-nearest
arithmetic, binary32 has a 24-bit significand and normal exponent range roughly
\(-126\) to \(+127\), while binary64 has a 53-bit significand and range
\(-1022\) to \(+1023\) \citep{ieee754_2019,goldberg_1991}. The corresponding
unit roundoff is about \(2^{-24}\) and \(2^{-53}\), respectively
\citep{higham_2002}; this is a local scale, not a solution-error bound for a
nonlinear time-dependent HRSC calculation. A unit in the last place (ULP) is
the spacing between adjacent representable numbers; later
\(\mathrm{ULP}_{\max}\) tables use it to expose bit-level differences hidden
by relative norms.

The key difficulty is non-associativity, with a related hazard of catastrophic
cancellation: subtraction of nearly-equal quantities loses leading digits and
inflates the relative error. In binary64, \((10^{-18}+1)-1\) rounds the small
increment away and returns zero, whereas \(10^{-18}+(1-1)\) returns
\(10^{-18}\). Such changes alter wave speeds, limiters, fluxes, and later steps. Compensated (Kahan) summation tracks lost low-order bits in a
correction term, recovering near full precision in long accumulations;
Algorithm~1 uses it for the simulation-time counter. The report measures the size and location of these effects rather than
assuming harmlessness.

Compiler and device choices add variation. Ordinary \texttt{-O3} may change
instruction selection, but fast-math or \texttt{-Ofast} can permit
reassociation, FMA contraction, reciprocal-based division, and approximate
square-root paths. Since floating-point addition is non-associative, GPU block
or warp scheduling and atomic reductions can change the accumulation tree; the
solver comparisons below avoid summation-order claims unless an ordered max/min
or saved-state metric supports them. Reproducible summation algorithms can
control this source of variation, but they must be designed explicitly rather
than assumed from the mathematical expression \citep{demmel_nguyen_2013}. CUDA
devices follow IEEE-754 and expose fused multiply-add (FMA), which aids
accuracy and throughput but shifts bit patterns relative to a non-FMA CPU
build \citep{whitehead_fitflorea_2011}. Comparison order can also matter for tied
or nearly tied candidates. CFD precision studies such as Brogi et al.'s
OpenFOAM work show why reduced or mixed precision must be validated case by
case \citep{brogi_etal_2024}. This report also uses Verificarlo, which
implements Monte Carlo arithmetic (MCA) by perturbing floating-point operations
to estimate rounding sensitivity \citep{parker_1997,denis_etal_2016}.

\paragraph{Virtual precision versus IEEE binary32.}
Verificarlo's \texttt{p8}, \texttt{p16}, \texttt{p32}, and \texttt{p53}
settings are virtual mantissa precisions used by the MCA-RR operator, which
draws an independent random rounding for every elementary floating-point
operation. An IEEE binary32 run, by contrast, stores every value with a
fixed 24-bit significand and rounds each operation deterministically under
round-to-nearest-even semantics \citep{ieee754_2019}. The two differ in
both the rounding rule and the noise model, so \texttt{p32} and IEEE fp32
are not interchangeable: \texttt{p32} probes rounding sensitivity at a
virtual mantissa width, whereas IEEE fp32 is a storage and arithmetic
format used in actual production runs. Wherever this report quantifies an
fp32 effect, the numbers come from a real \texttt{HRSC\_REAL=float} build,
not a \texttt{p32} sample.
% <<SECTION_4_END>>
